\assignq{3}{1}{Finite nuclear size of hydrogen-like \(^{209}\mathrm{Bi}\)}{1}
Work in atomic units. A point-nucleus Coulomb potential is
\(V_0=-Z/r\), while the specified uniform sphere has
\[
V_R(r)=
\begin{cases}
-\frac Z{2R}(3-r^2/R^2)&r<R,\\
-Z/r&r\ge R.
\end{cases}
\]
Thus the perturbation is
\[
\Delta V(r)=Z\left(\frac1r-\frac3{2R}
+\frac{r^2}{2R^3}\right)\quad(r<R),
\]
and zero for \(r\ge R\).

\subsection*{(a) Derive the leading \(s\)-state correction}
First-order perturbation theory yields
\[
\Delta E_{ns}^{(1)}
=\mel{ns}{\Delta V}{ns}
\simeq4\pi|\psi_{ns}(0)|^2\int_0^R
r^2\Delta V(r)\diff r .
\]
For hydrogenic \(s\)-states,
\(|\psi_{ns}(0)|^2=Z^3/(\pi n^3)\).
Perform the radial integral explicitly:
\[
\begin{aligned}
\int_0^R r^2\Delta V(r)\diff r
&=Z\left[\frac{R^2}{2}
-\frac{3R^2}{6}+\frac{R^2}{10}\right]\\
&=\frac{ZR^2}{10}.
\end{aligned}
\]
Hence
\[
\boxed{\Delta E_{ns}^{\rm fns}
=\frac{2Z^4R^2}{5n^3}
=\frac{2Z^2}{5n^3}(ZR)^2\quad(\text{Hartree})}.
\]
It is positive: spreading out the nuclear charge
weakens the strongly attractive potential near the origin.

\subsection*{(b) Numerical integrals for \(^{209}\mathrm{Bi}^{82+}\)}
The nuclear charge is \(Z=83\). With
\(R_{\rm rms}=5.5211\,\mathrm{fm}\), take
\[
R=\sqrt{\frac53}R_{\rm rms}=7.1277\,\mathrm{fm}
=1.34694\times10^{-4}a_0.
\]
Use normalised hydrogenic \(R_{nl}(r)\) and integrate
\[
\boxed{\Delta E_{nl}=\int_0^R
|R_{nl}(r)|^2r^2\Delta V(r)\diff r}.
\]
The accompanying Python solver performs this
quadrature rather than assuming that all \(l\) give
the same \(s\)-state correction:
\[
\begin{array}{c|r|r}
\text{State}&\Delta E_{\rm fns}\ (\mathrm{Ha})&
\Delta E_{\rm fns}\ (\mathrm{eV})\\\hline
1s&0.3412153&9.28494\\
2p&9.54585\times10^{-8}&2.59756\times10^{-6}\\
3d&8.63359\times10^{-15}&2.34932\times10^{-13}
\end{array}
\]
The leading-order \(1s\) formula gives
\(0.3444055\,\mathrm{Ha}\), approximately
\(0.94\%\) higher than numerical first-order
Schr\"odinger perturbation theory. The higher
angular momenta are strongly suppressed near
the nuclear region because \(R_{nl}\sim r^l\).

\subsection*{(c) Relativistic comparison}
The exact \emph{point-nucleus} Dirac Coulomb
energy (excluding rest mass) is
\[
E_{n\kappa}^{D}=\frac1{\alpha_{\rm fs}^2}
\left[
\left(1+\frac{(Z\alpha_{\rm fs})^2}
{(n-|\kappa|+\sqrt{\kappa^2-(Z\alpha_{\rm fs})^2})^2}
\right)^{-1/2}-1\right]\ {\rm Ha}.
\]
Compare with \(E_n^{NR}=-Z^2/(2n^2)\):
\[
\begin{array}{c|r|r}
\text{Fine-structure level}&
E^{D}-E^{NR}\ (\mathrm{Ha})&
\Delta E_{\rm fns}^{NR}\ (\mathrm{Ha})\\\hline
1s_{1/2}&-391.87&0.3412\\
2p_{1/2}&-123.80&9.55\times10^{-8}\\
3d_{3/2}&-12.31&8.63\times10^{-15}
\end{array}
\]
The relativistic shifts are far larger than
these \emph{nonrelativistic} finite-nucleus
perturbation results. This is not a precision
relativistic heavy-ion calculation:
\(Z\alpha_{\rm fs}\simeq0.606\), so the real
Dirac radial functions have significantly
different nuclear-region amplitudes.
\methodnote{Numerical source:
\texttt{scripts/assignment\_numerical\_solutions.py}.
For \(R_{nl}\) use \(r^2\diff r\); for
\(u=rR_{nl}\) use \(\diff r\).}
